\documentclass[12pt]{article}
\usepackage[utf8]{inputenc}
\usepackage[T1]{fontenc}
\usepackage[margin=1in]{geometry}
\usepackage{setspace}
\setlength{\parindent}{2em}
\setlength{\parskip}{1.25em}
\renewcommand{\baselinestretch}{1.0}
\makeatletter
\newcommand{\@subtitle}{}
\newcommand{\subtitle}[1]{\renewcommand{\@subtitle}{#1}}
\renewcommand{\@maketitle}{%
  \newpage
  \null
  \vskip 1em%
  \begin{center}%
    {\LARGE \@title \par}%
    \ifx\@subtitle\@empty\else
      \vskip 0.5em%
      {\large \@subtitle \par}%
    \fi
  \end{center}%
  \par
  \vskip 1em}
\makeatother
\usepackage{xcolor}
\usepackage[colorlinks=true,urlcolor=blue!55!black,linkcolor=black]{hyperref}

\title{\textit{Buckley v. Valeo}, 424 U.S. 1 (1976)}
\subtitle{Case Brief}
\author{}
\date{}

\begin{document}
\maketitle

\section{Memory Jogger}

In \href{https://www.courtlistener.com/opinion/109380/buckley-v-valeo/}{\textit{Buckley v. Valeo}}, the Supreme Court established the foundational framework of modern campaign finance law by upholding limits on direct political contributions, disclosure rules, and voluntary presidential public financing while striking down limits on campaign and independent expenditures under the First Amendment and invalidating congressional appointment of Federal Election Commission members under the Appointments Clause.

\section{Detailed Case Facts}

Following the Watergate scandal and revelations of widespread, undisclosed, and illegal campaign financing during the 1972 presidential election, Congress enacted sweeping 1974 amendments to the Federal Election Campaign Act of 1971 (FECA) and related provisions of the Internal Revenue Code of 1954 (\href{https://supreme.justia.com/cases/federal/us/424/1/}{424 U.S. at 6--7}; \href{https://www.cambridge.org/core/books/law-of-freedom/campaign-finance/67D70830F386B05225DCEF7D8D8C25C7}{Cambridge Univ. Press, \textit{The Law of Freedom}}). The D.C. Circuit described the statute as ``by far the most comprehensive reform legislation [ever] passed by Congress concerning the election of the President, Vice-President, and members of Congress'' (\href{https://supreme.justia.com/cases/federal/us/424/1/}{424 U.S. at 7}).

The challenged statutory regime contained four major pillars (\href{https://supreme.justia.com/cases/federal/us/424/1/}{424 U.S. at 1--2, 7}):
\begin{enumerate}
  \item \textbf{Contribution and Expenditure Limitations (18 U.S.C. \S 608):} The Act capped individual and group contributions at \$1,000 per candidate per election (\S 608(b)(1)), political committee contributions at \$5,000 per candidate per election (\S 608(b)(2)), and total annual contributions by an individual at \$25,000 (\S 608(b)(3)), while capping certain exempt volunteer incidental expenses at \$500 (\S 591(e)(5)). On the spending side, it capped independent expenditures ``relative to a clearly identified candidate'' at \$1,000 per calendar year (\S 608(e)(1)), capped a candidate's use of personal or immediate family funds at \$50,000 for President/Vice President, \$35,000 for Senate, and \$25,000 for House races (\S 608(a)(1)), and imposed overall primary and general election campaign spending ceilings (\S 608(c)).
  \item \textbf{Recordkeeping and Disclosure Requirements (2 U.S.C. \S\S 431--438):} Political committees had to record the name and address of every contributor giving over \$10 (\S 432(c)(2)) and publicly report the name, address, occupation, and principal place of business of every contributor giving over \$100 (\S 434(b)). Individuals and non-committee groups making independent contributions or expenditures exceeding \$100 per year ``for the purpose of ... influencing'' a federal election also had to file statements with the Commission (\S 434(e)).
  \item \textbf{Public Financing of Presidential Campaigns (Subtitle H of the IRC, 26 U.S.C. \S\S 6096, 9001--9042):} Financed via an optional \$1 (\$2 joint) income tax check-off, Subtitle H provided public grants for nominating conventions and general elections, dividing parties into ``major'' parties (25\% or more of the prior presidential vote, receiving full funding), ``minor'' parties (5\% to under 25\%, receiving proportional funding), and ``new'' parties (under 5\%, eligible only for post-election funds upon reaching 5\%), alongside matching funds for primary candidates who raised \$5,000 in each of 20 states in contributions of \$250 or less. Accepting public funds required agreeing to expenditure caps.
  \item \textbf{The Federal Election Commission (2 U.S.C. \S 437c):} The Act created an eight-member Federal Election Commission (FEC) with recordkeeping, investigative, rulemaking, advisory-opinion, adjudicatory, and civil enforcement powers, including authority to disqualify candidates who failed to file required reports. Two members were ex officio nonvoting congressional officers (the Secretary of the Senate and the Clerk of the House). Of the six voting members, two were appointed by the President pro tempore of the Senate, two by the Speaker of the House, and two by the President, with all six requiring confirmation by a majority of both Houses of Congress.
\end{enumerate}

On January 2, 1975, an ideologically diverse coalition of plaintiffs---including Conservative U.S. Senator James L. Buckley of New York, liberal former Senator and presidential candidate Eugene J. McCarthy of Minnesota, the Conservative Party of New York, the Mississippi Republican Party, the Libertarian Party, the New York Civil Liberties Union, the American Conservative Union, and \textit{Human Events, Inc.}---filed suit against Francis R. Valeo (Secretary of the Senate and ex officio FEC member), the Clerk of the House, the FEC, the Attorney General, and the Comptroller General (\href{https://supreme.justia.com/cases/federal/us/424/1/}{424 U.S. at 7--8}; \href{https://www.fec.gov/legal-resources/court-cases/buckley-v-valeo/}{FEC Case Summary}). Plaintiffs sought declaratory and injunctive relief under the First Amendment, the Fifth Amendment Due Process Clause (incorporating equal protection), the General Welfare Clause (Art. I, \S 8), and the Appointments Clause (Art. II, \S 2, cl. 2).

\section{Procedural History}

Plaintiffs filed their action in the U.S. District Court for the District of Columbia, invoking both general federal question jurisdiction (28 U.S.C. \S\S 1331, 2201--2202) and FECA's expedited judicial review mechanism (2 U.S.C. \S 437h(a)) (\href{https://supreme.justia.com/cases/federal/us/424/1/}{424 U.S. at 8--9}). The district judge initially denied plaintiffs' request for a three-judge district court and transmitted the case to the U.S. Court of Appeals for the D.C. Circuit. The D.C. Circuit, sitting en banc, remanded for factfinding and certification of constitutional questions.

On August 15, 1975, sitting en banc, the D.C. Circuit (\textit{Buckley v. Valeo}, 519 F.2d 821 (D.C. Cir. 1975)) upheld nearly all challenged FECA provisions---including contribution limits, expenditure limits, disclosure mandates, and the FEC's appointment structure---relying on \textit{United States v. O'Brien}, 391 U.S. 367 (1968), to treat money as regulable conduct rather than pure speech, and finding a compelling interest in preserving electoral integrity (\href{https://supreme.justia.com/cases/federal/us/424/1/}{424 U.S. at 10, 15--16}; \href{https://www.fec.gov/legal-resources/court-cases/buckley-v-valeo/}{FEC Case Summary}). Separately, a three-judge district court (\textit{Buckley v. Valeo}, 401 F. Supp. 1235 (D.D.C. 1975)) upheld the Subtitle H public financing provisions (\href{https://supreme.justia.com/cases/federal/us/424/1/}{424 U.S. at 2, 86}).

Plaintiffs appealed both judgments directly to the U.S. Supreme Court (No. 75-436 from the D.C. Circuit; No. 75-437 from the three-judge district court). Oral argument was held on November 10, 1975 (\href{https://www.oyez.org/cases/1975/75-436}{Oyez}). On January 30, 1976, the Supreme Court issued a lengthy \textit{per curiam} opinion affirming in part and reversing in part the D.C. Circuit (No. 75-436) and affirming the three-judge district court (No. 75-437) (\href{https://supreme.justia.com/cases/federal/us/424/1/}{424 U.S. at 5}).

\section{Judicial Votes}

Eight Justices participated. Justice William O. Douglas did not participate at oral argument and retired two days later on November 12, 1975; his successor, Justice John Paul Stevens, took no part in the consideration or decision of the cases (\href{https://supreme.justia.com/cases/federal/us/424/1/}{424 U.S. at 5}; \href{https://ballotpedia.org/Buckley_v._Valeo}{Ballotpedia}).

\begin{itemize}
  \item \textbf{Per Curiam Opinion (Majority):} All eight participating Justices joined the threshold Article III ``case or controversy'' holding (424 U.S. at 11--12). Justices William J. Brennan, Jr., Potter Stewart, and Lewis F. Powell, Jr., joined the \textit{per curiam} opinion in full. Five other Justices joined specific parts while filing separate opinions concurring in part and dissenting in part, yielding the following issue-by-issue majorities (\href{https://supreme.justia.com/cases/federal/us/424/1/}{424 U.S. at 5}):
  \begin{itemize}
    \item \textit{Part I-B (Contribution limits upheld):} 6--2 (Brennan, Stewart, Powell, Marshall, Rehnquist, and White; Burger and Blackmun dissenting).
    \item \textit{Part I-C-1 \& I-C-3 (Independent and overall campaign expenditure limits struck down):} 7--1 (Brennan, Stewart, Powell, Marshall, Blackmun, Rehnquist, and Burger; White dissenting).
    \item \textit{Part I-C-2 (Personal/family funds expenditure limit struck down):} 6--2 (Brennan, Stewart, Powell, Blackmun, Rehnquist, and Burger; Marshall and White dissenting).
    \item \textit{Part II (Recordkeeping and disclosure requirements upheld):} 7--1 (Brennan, Stewart, Powell, Marshall, Blackmun, Rehnquist, and White; Burger dissenting as to small-contribution disclosure).
    \item \textit{Part III (Subtitle H public financing upheld):} 6--2 on general election financing (Part III-B-1) (Brennan, Stewart, Powell, Marshall, Blackmun, and White; Burger and Rehnquist dissenting); 7--1 on the remainder of Part III (Burger dissenting).
    \item \textit{Part IV (FEC appointment structure struck down under Appointments Clause):} 7--1 (Brennan, Stewart, Powell, Marshall, Blackmun, Rehnquist, and Burger [except Burger dissented from according \textit{de facto} validity to past FEC acts]; White dissenting).
  \end{itemize}
  \item \textbf{Concurrences and Dissents in Part:} Chief Justice Warren E. Burger (424 U.S. at 235), Justice Byron R. White (424 U.S. at 257), Justice Thurgood Marshall (424 U.S. at 286), Justice Harry A. Blackmun (424 U.S. at 290), and Justice William H. Rehnquist (424 U.S. at 290) each filed opinions concurring in part and dissenting in part (\href{https://supreme.justia.com/cases/federal/us/424/1/}{424 U.S. at 5}).
  \item \textbf{Not Participating:} Justice John Paul Stevens (and Justice William O. Douglas, retired).
\end{itemize}

\section{Holding}

\textbf{Judgment:} Affirmed in part and reversed in part as to the D.C. Circuit (No. 75-436); affirmed as to the three-judge District Court (No. 75-437). The Court accorded \textit{de facto} validity to past acts of the FEC and stayed its judgment on Part IV for 30 days to allow Congress to reconstitute the Commission (\href{https://supreme.justia.com/cases/federal/us/424/1/}{424 U.S. at 5}).

\textbf{Rules Announced:}
\begin{enumerate}
  \item \textbf{Contribution vs. Expenditure Distinction:} Political spending is directly tied to First Amendment expression and is not mere conduct subject to intermediate scrutiny under \textit{United States v. O'Brien}. However, contribution limits entail only a ``marginal restriction'' on speech and are constitutional if ``closely drawn'' to match the ``sufficiently important interest'' of preventing actual \textit{quid pro quo} corruption or the appearance of corruption (\href{https://supreme.justia.com/cases/federal/us/424/1/}{424 U.S. at 20--21, 25--27}). By contrast, expenditure limits---whether on independent speakers, candidates' personal funds, or overall campaign spending---impose direct and substantial restraints on the quantity of core political speech, trigger ``exacting scrutiny,'' and cannot be justified by either anti-corruption or political-equalization interests (\href{https://supreme.justia.com/cases/federal/us/424/1/}{424 U.S. at 39--59}).
  \item \textbf{Express Advocacy Narrowing Construction:} To avoid unconstitutional vagueness under the First Amendment, independent expenditure rules (and reporting under \S 434(e)) must be narrowly construed to reach only communications containing explicit words advocating the election or defeat of a clearly identified candidate---such as ``vote for,'' ``elect,'' ``support,'' or ``cast your ballot for'' (\href{https://supreme.justia.com/cases/federal/us/424/1/}{424 U.S. at 44 \& n.52, 80}; \href{https://www.fec.gov/resources/legal-resources/litigation/Buckley.pdf}{FEC Slip Op.}).
  \item \textbf{Disclosure and Minor-Party Exemptions:} Mandatory recordkeeping (\$10 threshold) and public disclosure (\$100 threshold) survive exacting scrutiny because they serve substantial governmental interests in informing voters, deterring corruption, and enforcing contribution limits. Minor parties are not entitled to a blanket exemption, but may obtain an as-applied exemption upon showing a ``reasonable probability'' that disclosure will expose contributors to threats, harassment, or reprisals (\href{https://supreme.justia.com/cases/federal/us/424/1/}{424 U.S. at 66--74, 82--84}).
  \item \textbf{Voluntary Public Financing:} Congress may use the General Welfare Clause (Art. I, \S 8) to fund presidential campaigns and may condition acceptance of public funds on voluntary agreement to expenditure caps; differentiating between major, minor, and new parties based on prior vote totals does not invidiously discriminate under the Fifth Amendment (\href{https://supreme.justia.com/cases/federal/us/424/1/}{424 U.S. at 90--109}).
  \item \textbf{Appointments Clause:} Any appointee ``exercising significant authority pursuant to the laws of the United States'' is an ``Officer of the United States'' who must be appointed pursuant to Article II, \S 2, cl. 2. Because four of the FEC's six voting members were appointed by congressional leaders, the FEC could not constitutionally exercise rulemaking, adjudicatory, or civil enforcement powers, though it could exercise purely investigative and informative functions (\href{https://supreme.justia.com/cases/federal/us/424/1/}{424 U.S. at 126, 137--141}).
\end{enumerate}

\textbf{Holding vs. Dicta:} The Court's invalidation of \S\S 608(a), 608(c), and 608(e)(1), its upholding of \S 608(b) contribution caps, \S 434 disclosure provisions, and Subtitle H public financing, and its Appointments Clause ruling on the FEC are binding holdings. By contrast, footnote 52's illustrative list of ``magic words'' of express advocacy was explanatory guidance that later courts treated as a bright-line test until modified in \textit{McConnell v. FEC} (2003), and the Court's discussion of when future minor-party harassment showings would warrant an as-applied disclosure exemption was forward-looking dicta later applied in \textit{Brown v. Socialist Workers '74 Campaign Committee} (1982).

\section{Analysis of Opinions}

\subsection{Per Curiam Opinion (Majority)}

The \textit{per curiam} opinion---jointly drafted by several Justices to expedite decision before the 1976 election---rejected the D.C. Circuit's reliance on \textit{United States v. O'Brien} and time, place, and manner doctrine. Because ``virtually every means of communicating ideas in today's mass society requires the expenditure of money,'' restricting political spending directly restricts the quantity and diversity of political speech (\href{https://supreme.justia.com/cases/federal/us/424/1/}{424 U.S. at 16--19}).

Crucially, however, the Court bifurcated campaign finance into contributions and expenditures. A contribution is a ``general expression of support'' whose communicative value ``does not increase perceptibly with the size of his contribution,'' and converting a contribution into speech requires action by someone other than the donor (\href{https://supreme.justia.com/cases/federal/us/424/1/}{424 U.S. at 20--21}). Because large contributions carry a tangible risk of \textit{quid pro quo} corruption and erode public confidence through the appearance of impropriety, the \$1,000, \$5,000, and \$25,000 caps were closely drawn to a weighty interest, and courts lack a ``scalpel'' to second-guess whether \$2,000 would work as well as \$1,000 (\href{https://supreme.justia.com/cases/federal/us/424/1/}{424 U.S. at 26--30, 38}).

Expenditure ceilings failed exacting scrutiny. Independent expenditures lacking prearrangement or coordination with a candidate pose a ``substantially diminished potential for abuse,'' and narrowing \S 608(e)(1) to express advocacy to cure vagueness rendered it an ineffective loophole-closer anyway (\href{https://supreme.justia.com/cases/federal/us/424/1/}{424 U.S. at 44--47}). Candidate self-funding actually \textit{reduces} dependence on outside donors (\href{https://supreme.justia.com/cases/federal/us/424/1/}{424 U.S. at 53}). Most famously, the Court flatly rejected the equalization rationale: ``the concept that government may restrict the speech of some elements of our society in order to enhance the relative voice of others is wholly foreign to the First Amendment'' (\href{https://supreme.justia.com/cases/federal/us/424/1/}{424 U.S. at 48--49}).

On disclosure, the Court applied \textit{NAACP v. Alabama} exacting scrutiny but found three substantial interests---voter information, corruption deterrence (citing Justice Brandeis's maxim that ``[s]unlight is said to be the best of disinfectants''), and enforcement of contribution caps---sufficient to justify the \$10/\$100 thresholds absent record proof of harassment (\href{https://supreme.justia.com/cases/federal/us/424/1/}{424 U.S. at 64--68, 82--84}). On Subtitle H, the Court held that public financing enlarges rather than abridges speech and that requiring a ``significant modicum of support'' avoids funding frivolous candidacies and splintered factionalism (\href{https://supreme.justia.com/cases/federal/us/424/1/}{424 U.S. at 92--96}). Finally, on the FEC, the Court invoked Madison's \textit{Federalist Nos. 47 and 51} and the drafting history of the Constitutional Convention to hold that Congress cannot aggrandize itself by appointing ``Officers of the United States'' who execute federal law (\href{https://supreme.justia.com/cases/federal/us/424/1/}{424 U.S. at 120--132}).

\subsection{Chief Justice Burger (Concurring in Part and Dissenting in Part)}

Chief Justice Burger joined Parts I-C (invalidating expenditure caps) and IV (invalidating the FEC's appointment structure, except he dissented from granting \textit{de facto} validity to past FEC actions), but dissented from upholding contribution limits, small-contribution disclosure, and Subtitle H public financing (\href{https://supreme.justia.com/cases/federal/us/424/1/}{424 U.S. at 5, 235}; \href{https://ballotpedia.org/Buckley_v._Valeo}{Ballotpedia}). Rejecting the majority's core distinction, Burger famously argued that ``contributions and expenditures are two sides of the same First Amendment coin'' (\href{https://edge.sagepub.com/epsteinshort9e/student-resources/chapter-5-the-executive/buckley-v-valeo-1976}{424 U.S. at 241}). He viewed the \$10 and \$100 disclosure thresholds as intrusive burdens on grassroots political privacy and warned that government financing of campaigns invited unconstitutional governmental entanglement in the political process.

\subsection{Justice White (Concurring in Part and Dissenting in Part)}

Justice White would have upheld virtually the entire statute---joining Part III (public financing) and concurring in upholding contribution limits and disclosure, while dissenting from the invalidation of expenditure ceilings and the FEC's appointment structure (\href{https://supreme.justia.com/cases/federal/us/424/1/}{424 U.S. at 5, 257}). White contended that ``money is not speech'' (\href{https://edge.sagepub.com/epsteinshort9e/student-resources/chapter-5-the-executive/buckley-v-valeo-1976}{424 U.S. at 262--263}) and argued that Congress should receive substantial deference in determining that expenditure ceilings were integral to preventing corruption, curbing skyrocketing campaign costs, and maintaining a workable regulatory scheme. On Article II, he argued that Congress's plenary authority over federal elections permitted legislative participation in appointing FEC commissioners.

\subsection{Justice Marshall (Concurring in Part and Dissenting in Part)}

Justice Marshall joined the \textit{per curiam} opinion in every respect except Part I-C-2, dissenting solely from the invalidation of \S 608(a)'s cap on a candidate's personal and family expenditures (\href{https://supreme.justia.com/cases/federal/us/424/1/}{424 U.S. at 5, 286}; \href{https://ballotpedia.org/Buckley_v._Valeo}{Ballotpedia}). Marshall argued that once the Court upheld limits on outside contributions, permitting wealthy candidates to spend unlimited personal fortunes created a severe, asymmetric advantage for the rich and undermined the equal-opportunity interest in allowing non-wealthy citizens to compete for public office.

\subsection{Justice Blackmun (Concurring in Part and Dissenting in Part)}

Justice Blackmun joined all of the \textit{per curiam} opinion except Part I-B, dissenting from the Court's decision to uphold contribution limitations (\href{https://supreme.justia.com/cases/federal/us/424/1/}{424 U.S. at 5, 290}). Like Chief Justice Burger, Blackmun saw no principled constitutional distinction between political contributions and political expenditures for First Amendment purposes and would have invalidated both.

\subsection{Justice Rehnquist (Concurring in Part and Dissenting in Part)}

Justice Rehnquist joined the entire \textit{per curiam} opinion except Part III-B-1, dissenting from the Court's approval of Subtitle H's general-election public financing formula as applied to minor and new political parties (\href{https://supreme.justia.com/cases/federal/us/424/1/}{424 U.S. at 5, 290}). Rehnquist argued that Congress had impermissibly enshrined the two existing major parties while denying pre-election funding to minor and new parties, violating the equal protection principles embedded in the Fifth Amendment.

\section{Examples: Future Applications}

\begin{enumerate}
  \item \textbf{Same-Side Application 1 (State Contribution Caps):} A state enacts a \$1,075 per-election cap on individual contributions to candidates for statewide office to combat actual and perceived corruption. Under \textit{Buckley}, this law is \textbf{constitutional}. Because contribution limits impose only a marginal burden on speech, the state need only show a sufficiently important anti-corruption interest and closely drawn means, without courts acting as a ``scalpel'' on the exact dollar amount (see \href{https://www.oyez.org/cases/1999/98-963}{\textit{Nixon v. Shrink Missouri Government PAC}, 528 U.S. 377 (2000)}).
  \item \textbf{Same-Side Application 2 (Independent Expenditure Ban by Outside Groups):} Congress passes a statute capping independent political action committees (PACs) at \$5,000 in uncoordinated television ads expressly advocating the defeat of a presidential candidate. Under \textit{Buckley}, this cap is \textbf{unconstitutional}. Independent expenditures uncoordinated with a candidate constitute core political expression subject to exacting scrutiny, and neither anti-corruption nor equalization interests justify capping them (see \href{https://supreme.justia.com/cases/federal/us/470/480/}{\textit{FEC v. NCPAC}, 470 U.S. 480 (1985)}).
  \item \textbf{Opposite-Side Application 1 (Aggregate Biennial Contribution Limits):} Congress imposes an aggregate ceiling limiting the total amount a single donor may contribute across \textit{all} federal candidates and party committees combined in an election cycle, even when every individual donation stays within the base per-candidate limit. Although \textit{Buckley} briefly upheld FECA's \$25,000 aggregate cap as an anti-circumvention measure (424 U.S. at 38), applying \textit{Buckley}'s strict \textit{quid pro quo} corruption framework today results in \textbf{invalidation} where aggregate limits restrict how many candidates a citizen may support without targeting actual circumvention (see \href{https://www.oyez.org/cases/2013/12-536}{\textit{McCutcheon v. FEC}, 572 U.S. 185 (2014)}).
  \item \textbf{Opposite-Side Application 2 (As-Applied Disclosure Exemption for Harassed Minor Party):} A small ideological party with a documented history of FBI surveillance, employer firings, and physical threats against its known supporters challenges a state law requiring public disclosure of all donors giving over \$100. Unlike the plaintiffs in \textit{Buckley}, who offered only speculative fears, this party satisfies \textit{Buckley}'s ``reasonable probability'' standard of threats, harassment, or reprisals (\href{https://supreme.justia.com/cases/federal/us/424/1/}{424 U.S. at 74}) and is \textbf{constitutionally exempt} from disclosure.
  \item \textbf{Fence-Sitter (Trigger-Based Public Financing Matching Funds):} A state creates a voluntary public financing system (permitted under \textit{Buckley}) in which participating candidates receive a fixed grant plus extra ``matching funds'' dollar-for-dollar whenever a privately funded opponent or independent group spends above a threshold. On one hand, the program is voluntary and adds rather than caps speech; on the other, triggering extra government subsidies for one's opponent directly penalizes a nonparticipating candidate's or independent speaker's protected expenditures. This sits squarely on the constitutional fault line between \textit{Buckley}'s approval of voluntary public financing and its ban on penalizing political expenditures (ultimately struck down in \textit{Arizona Free Enterprise Club's Freedom Club PAC v. Bennett}, 564 U.S. 721 (2011)).
\end{enumerate}

\section{Critique}

\textit{Buckley v. Valeo} has long been described as a ``no-man's-land of contestation'' assailed from every ideological direction (\href{https://www.cambridge.org/core/books/law-of-freedom/campaign-finance/67D70830F386B05225DCEF7D8D8C25C7}{Cambridge Univ. Press, \textit{The Law of Freedom}}). Logically, critics across the spectrum argue that \textit{Buckley}'s central dichotomy between contributions and expenditures is doctrinally unstable. By upholding limits on the supply of campaign funds (contributions) while striking down limits on demand (expenditures), the Court left candidates facing unlimited spending pressure from wealthy self-funders and outside groups while forced to raise money in small increments---a hydraulic distortion that channeled billions of dollars into independent ``super PACs'' and undisclosed ``dark money'' groups (\href{https://www.americanprogress.org/article/the-corporate-power-reset-that-makes-citizens-united-irrelevant/}{Center for American Progress}). Moreover, the Court's ``magic words'' narrowing construction in footnote 52 proved notoriously easy to evade through ``issue ads'' that praised or attacked candidates without uttering ``vote for'' or ``defeat.''

\textbf{Progressive Critique:} Progressive scholars and reformers argue that \textit{Buckley} committed a profound category error by equating money with speech and by narrowing the permissible governmental interest to \textit{quid pro quo} bribery. By declaring political equalization ``wholly foreign to the First Amendment,'' \textit{Buckley} disabled Congress from addressing structural distortions in which vast private wealth drowns out ordinary citizens' voices and warps legislative agendas (\href{https://www.cambridge.org/core/books/law-of-freedom/campaign-finance/67D70830F386B05225DCEF7D8D8C25C7}{Cambridge Univ. Press}; \href{https://uknowledge.uky.edu/cgi/viewcontent.cgi?article=5550&context=klj}{Kentucky Law Journal}). Progressives emphasize Justice White's and Justice Marshall's warnings that democratic self-government requires protecting electoral equality, not merely policing explicit bribes.

\textbf{Originalist and Textualist Critique:} Originalist and libertarian jurists---most notably Justice Clarence Thomas in \href{https://www.oyez.org/cases/1999/98-963}{\textit{Nixon v. Shrink Missouri Government PAC}} and \href{https://www.oyez.org/cases/2013/12-536}{\textit{McCutcheon v. FEC}}---attack \textit{Buckley} from the opposite flank. Echoing Chief Justice Burger and Justice Blackmun, they contend that the text and original public meaning of the First Amendment draw no distinction between spending money to speak directly and pooling money with a candidate to amplify political advocacy. From this vantage point, \textit{Buckley}'s reduced scrutiny for contribution limits is an unprincipled judicial compromise that allows incumbent legislators to write self-serving fundraising caps restricting challengers (\href{https://www.cato.org/commentary/buckley-v-valeo-not-what-ails-american-democracy}{Cato Institute}). Conversely, originalists widely praise Part IV of \textit{Buckley} as a landmark restoration of formalist separation-of-powers and Appointments Clause constraints on the administrative state.

\section{Key Quotations}

\begin{quote}
A restriction on the amount of money a person or group can spend on political communication during a campaign necessarily reduces the quantity of expression by restricting the number of issues discussed, the depth of their exploration, and the size of the audience reached. This is because virtually every means of communicating ideas in today's mass society requires the expenditure of money.
\end{quote}
\noindent --- Per Curiam, \href{https://supreme.justia.com/cases/federal/us/424/1/}{424 U.S. at 19}.

\begin{quote}
By contrast with a limitation upon expenditures for political expression, a limitation upon the amount that any one person or group may contribute to a candidate or political committee entails only a marginal restriction upon the contributor's ability to engage in free communication.
\end{quote}
\noindent --- Per Curiam, \href{https://supreme.justia.com/cases/federal/us/424/1/}{424 U.S. at 20--21}.

\begin{quote}
But the concept that government may restrict the speech of some elements of our society in order to enhance the relative voice of others is wholly foreign to the First Amendment, which was designed ``to secure `the widest possible dissemination of information from diverse and antagonistic sources,''' and ```to assure unfettered interchange of ideas for the bringing about of political and social changes desired by the people.'''
\end{quote}
\noindent --- Per Curiam, \href{https://supreme.justia.com/cases/federal/us/424/1/}{424 U.S. at 48--49}.

\begin{quote}
In the free society ordained by our Constitution, it is not the government, but the people -- individually, as citizens and candidates, and collectively, as associations and political committees -- who must retain control over the quantity and range of debate on public issues in a political campaign.
\end{quote}
\noindent --- Per Curiam, \href{https://supreme.justia.com/cases/federal/us/424/1/}{424 U.S. at 57}.

\begin{quote}
We think its fair import is that any appointee exercising significant authority pursuant to the laws of the United States is an ``Officer of the United States,'' and must, therefore, be appointed in the manner prescribed by \S 2, cl. 2, of that Article.
\end{quote}
\noindent --- Per Curiam, \href{https://supreme.justia.com/cases/federal/us/424/1/}{424 U.S. at 126}.

\end{document}
